\clearpage
\begin{examplebox}
\textbf{例1｜同一个式子，为什么会有不同结论？（最关键）}
\begin{center}\begin{tikzpicture}[font=\small,>=stealth,line join=round]
\draw[gray!40,rounded corners=3pt] (0,0) rectangle (14.2,8.8);
\node at (7.1,8.200000000000001) {例1：同一个平方公式，输入范围不同，结论不同};
\begin{scope}[shift={(3.5,2.6)},x=1.05cm,y=1.05cm]\draw[->,gray] (-2.2,0)--(2.4000000000000004,0) node[right]{$x$};\draw[->,gray] (0,-0.3)--(0,4.7) node[left]{$y$};\begin{scope}\clip (-2.2,-0.3) rectangle (2.2,4.5);\draw[blue,thick,domain=-2.1:2.1,samples=120] plot(\x,{\x*\x});\end{scope}\fill[orange!85!black] (1,1) circle(2pt) node[above right]{$(1,1)$};\fill[orange!85!black] (-1,1) circle(2pt) node[above left]{$(-1,1)$};\draw[gray,dotted] (-1,1)--(1,1);\end{scope}\node at (3.5,1.45) {$D=\mathbb R$：是偶函数};\begin{scope}[shift={(10.5,2.6)},x=1.05cm,y=1.05cm]\draw[->,gray] (-2.2,0)--(2.4000000000000004,0) node[right]{$x$};\draw[->,gray] (0,-0.3)--(0,4.7) node[left]{$y$};\begin{scope}\clip (-2.2,-0.3) rectangle (2.2,4.5);\draw[blue,thick,domain=0:2.1,samples=120] plot(\x,{\x*\x});\end{scope}\fill[orange!85!black] (1,1) circle(2pt) node[above right]{$(1,1)$};\draw[red!70!black,thick] (-1-.08,1-.08)--(-1+.08,1+.08) (-1-.08,1+.08)--(-1+.08,1-.08);\node[above left,red!70!black] at (-1,1) {缺失};\end{scope}\node at (10.5,1.45) {$D=[0,+\infty)$：非奇非偶};
\node[text width=12.8cm,align=center] at (7.1,.65) {红叉表示该输入没有函数值，不是函数图像上的点。};
\end{tikzpicture}\end{center}



判断$f(x)=x^2$在下列两个定义域上的奇偶性：\textbf{（1）}$D=\mathbb R$；\textbf{（2）}$D=[0,+\infty)$。

\textbf{（1）}定义域是全体实数，$x$与$-x$都可以输入；$f(-x)=(-x)^2=x^2=f(x)$。\textbf{是偶函数。}

\textbf{（2）}在$D=[0,+\infty)$中，$2$可以输入，但$-2$不可以输入，定义域不对称。\textbf{既不是奇函数，也不是偶函数。}

\textbf{本题收获：}表达式长得一样，不代表函数性质一定一样，\textbf{定义域也属于函数的一部分}。

\medskip
\end{examplebox}
\clearpage
\begin{examplebox}
\textbf{例2｜常规判断奇函数}
\begin{center}\begin{tikzpicture}[font=\small,>=stealth,line join=round]
\draw[gray!40,rounded corners=3pt] (0,0) rectangle (14.2,8.5);
\node at (7.1,7.9) {例2：不必一直递增，奇函数只要求对称伙伴};
\begin{scope}[shift={(3.6,4.2)},x=1.15cm,y=0.85cm]\draw[->,gray] (-2.2,0)--(2.4000000000000004,0) node[right]{$x$};\draw[->,gray] (0,-3)--(0,3.2) node[left]{$y$};\begin{scope}\clip (-2.2,-3) rectangle (2.2,3);\draw[blue,thick,domain=-2.2:2.2,samples=120] plot(\x,{\x*\x*\x-2*\x});\end{scope}\draw[orange!85!black,dashed] (-1,1)--(1,-1);\fill[orange!85!black] (-1,1) circle(2pt) node[above left]{$(-1,1)$};\fill[orange!85!black] (1,-1) circle(2pt) node[below right]{$(1,-1)$};\end{scope}\node[align=left,anchor=west] at (7.8,4.7) {$f(x)=x^3-2x$\\[12pt]定义域 $\mathbb R$ 成对\\[12pt]$f(-x)=-f(x)$\\[12pt]任意点绕原点半转后仍在曲线上};
\node[text width=12.8cm,align=center] at (7.1,.65) {橙色点对只是示例；正文的恒等式才证明对每个输入都成立。};
\end{tikzpicture}\end{center}



判断$f(x)=x^3-2x$的奇偶性。

\textbf{第1步}：定义域是$\mathbb R$，关于原点对称。

\textbf{第2步}：
\[
f(-x)=(-x)^3-2(-x)=-x^3+2x=-(x^3-2x)=-f(x).
\]
\textbf{结论：}是\textbf{奇函数}。

\medskip
\end{examplebox}
\clearpage
\begin{examplebox}
\textbf{例3｜没有$x=0$，仍然可以是奇函数}
\begin{center}\begin{tikzpicture}[font=\small,>=stealth,line join=round]
\draw[gray!40,rounded corners=3pt] (0,0) rectangle (14.2,9.4);
\node at (7.1,8.8) {例3：原点不在定义域，也可以关于原点对称};
\begin{scope}[shift={(3.6,4.9)},x=1.05cm,y=0.6cm]\draw[->,gray] (-2.6,0)--(2.8000000000000003,0) node[right]{$x$};\draw[->,gray] (0,-5.5)--(0,5.7) node[left]{$y$};\begin{scope}\clip (-2.6,-5.5) rectangle (2.6,5.5);\draw[blue,thick,domain=-2.6:-0.185,samples=120] plot(\x,{\x+1/\x});\draw[blue,thick,domain=0.185:2.6,samples=120] plot(\x,{\x+1/\x});\end{scope}\draw[orange!85!black,dashed] (-1,-2)--(1,2);\fill[orange!85!black] (-1,-2) circle(2pt) node[above left]{$(-1,-2)$};\fill[orange!85!black] (1,2) circle(2pt) node[below right]{$(1,2)$};\draw[red!65!black] (0,0) circle(2pt);\node[below right,red!65!black] at (0,0) {原点不在图上};\end{scope}\node[align=left,anchor=west] at (7.9,5.3) {$f(x)=x+1/x$\\[12pt]$D=\mathbb R\setminus\{0\}$ 对称\\[12pt]$f(-x)=-f(x)$\\[12pt]奇函数可以没有 $x=0$};
\node[text width=12.8cm,align=center] at (7.1,.65) {红圈只提示原点缺失；曲线在 $x=0$ 两侧趋向无穷，视窗裁切不是端点。};
\end{tikzpicture}\end{center}



判断$f(x)=\dfrac{x^2+1}{x}$的奇偶性。

定义域为$x\ne0$；若$x\ne0$，则$-x\ne0$，定义域对称。接着计算：
\[
f(-x)=\frac{(-x)^2+1}{-x}=-\frac{x^2+1}{x}=-f(x).
\]
\textbf{结论：}是\textbf{奇函数}。\textbf{提醒：}奇函数并不要求$0$一定在定义域中。

\medskip
\end{examplebox}
\clearpage
\begin{examplebox}
\textbf{例4｜定义域对称，也可以既非奇也非偶}
\begin{center}\begin{tikzpicture}[font=\small,>=stealth,line join=round]
\draw[gray!40,rounded corners=3pt] (0,0) rectangle (14.2,9.0);
\node at (7.1,8.4) {例4：输入能成对，不等于输出符合奇偶规律};
\begin{scope}[shift={(3.6,2.9)},x=1.15cm,y=1.0cm]\draw[->,gray] (-2.2,0)--(2.4000000000000004,0) node[right]{$x$};\draw[->,gray] (0,-0.6)--(0,4.7) node[left]{$y$};\begin{scope}\clip (-2.2,-0.6) rectangle (2.2,4.5);\draw[blue,thick,domain=-2.2:2.2,samples=120] plot(\x,{\x*\x+\x});\end{scope}\fill[orange!85!black] (1,2) circle(2pt) node[right]{$(1,2)$};\fill[orange!85!black] (-1,0) circle(2pt) node[below left]{$(-1,0)$};\draw[orange!85!black,dotted] (-1,0)--(-1,2)--(1,2);\end{scope}\node[align=left,anchor=west] at (7.6,4.8) {$f(x)=x^2+x$\\[12pt]$D=\mathbb R$ 对称\\[12pt]$f(1)=2,\ f(-1)=0$\\[12pt]两值既不相等，也不互为相反数\\[12pt]$f(0)=0$ 也不能说明它是奇函数};
\node[text width=12.8cm,align=center] at (7.1,.65) {找到一对违例点，就能否定全称结论。};
\end{tikzpicture}\end{center}



判断$f(x)=x^2+x$的奇偶性。

定义域是$\mathbb R$，但$f(-x)=x^2-x$。这个式子既不总是等于$x^2+x$，也不总是等于$-x^2-x$。\textbf{结论：既不是奇函数，也不是偶函数。}

\textbf{快速否定：}取$x=1$，$f(1)=2$、$f(-1)=0$，既不相等，也不互为相反数，足以说明\textbf{两种都不满足}。注意：取一个数\textbf{可以推翻}一种性质，但取一个数符合要求\textbf{不能证明}一种性质。

\medskip
\end{examplebox}
\clearpage
\begin{examplebox}
\textbf{例5｜已知一个函数值，直接求另一边}
\begin{center}\begin{tikzpicture}[font=\small,>=stealth,line join=round]
\draw[gray!40,rounded corners=3pt] (0,0) rectangle (14.2,9.0);
\node at (7.1,8.4) {例5：上下离中线同为17，未知函数值随即确定};
\begin{scope}[shift={(3.8,5.6)},x=.7cm,y=.16cm]\draw[->,gray] (-4,0)--(4.1,0) node[right]{$x$};\draw[->,gray] (0,-27)--(0,13) node[left]{$y$};\draw[teal,dashed] (-4,-7)--(4,-7);\fill[blue] (3,10) circle(2pt) node[above left]{$(3,10)$};\fill[orange!85!black] (-3,-24) circle(2pt) node[below right]{$(-3,-24)$};\draw[gray,dotted] (3,10)--(-3,-24);\fill[teal] (0,-7) circle(2pt) node[below right]{$(0,-7)$};\draw[<->,blue] (3.6,-7)--(3.6,10) node[midway,right]{$17$};\draw[<->,orange!85!black] (-3.6,-24)--(-3.6,-7) node[midway,left]{$17$};\end{scope}\node[align=left,anchor=west] at (7.8,4.5) {$g(x)=f(x)+7$ 为奇函数\\[12pt]$g(3)=17,\ g(-3)=-17$\\[12pt]$f(3)+f(-3)=-14$\\[12pt]$f(-3)=-24$};
\node[text width=12.8cm,align=center] at (7.1,.65) {只画已知与必然的点及辅助线，不假设参数 $a,b$ 决定的具体曲线。};
\end{tikzpicture}\end{center}



已知$f(x)=x^5+ax^3+bx-7$，且$f(3)=10$，求$f(-3)$。

把末尾的$-7$暂时拿开，令$g(x)=f(x)+7=x^5+ax^3+bx$。因为$g(-x)=-g(x)$，$g$是奇函数。
\[
g(3)=f(3)+7=17,\quad g(-3)=-17.
\]
因此$f(-3)=g(-3)-7=\boxed{-24}$。

\textbf{更快的记法：}$f(-x)+7=-[f(x)+7]$，也就是$f(-x)+f(x)=-14$。代$x=3$也可立即得到答案。

\medskip
\end{examplebox}
\clearpage
\begin{examplebox}
\textbf{例6｜带参数：要让函数是偶函数，该怎样选$a$？}
\begin{center}\begin{tikzpicture}[font=\small,>=stealth,line join=round]
\draw[gray!40,rounded corners=3pt] (0,0) rectangle (14.2,9.0);
\node at (7.1,8.4) {例6：去掉奇的部分，左右才会等高};
\begin{scope}[shift={(3.6,3.3)},x=1.05cm,y=1.65cm]\draw[->,gray] (-2.5,0)--(2.7,0) node[right]{$x$};\draw[->,gray] (0,-0.2)--(0,2.5) node[left]{$y$};\begin{scope}\clip (-2.5,-0.2) rectangle (2.5,2.3);\draw[teal!80!black,thick,domain=-2.5:2.5,samples=120] plot(\x,{1});\draw[blue,thick,domain=-2.5:2.5,samples=120] plot(\x,{1+\x/(\x*\x+1)});\end{scope}\fill[blue] (-1,.5) circle(2pt) node[below left]{$(-1,1/2)$};\fill[blue] (1,1.5) circle(2pt) node[above right]{$(1,3/2)$};\end{scope}\node[align=left,anchor=west] at (7.8,4.8) {$f(x)=1+\dfrac{ax}{x^2+1}$\\[14pt]\textcolor{blue}{$a=1$：两边值不同}\\[12pt]\textcolor{teal!80!black}{$a=0$：恒为 $1$，是偶函数}\\[12pt]要对所有 $x$ 都相同，只能 $a=0$};
\node[text width=12.8cm,align=center] at (7.1,.65) {绿色恒等于1；蓝色为参数 $a=1$ 的反例。纵轴显示实际函数值。};
\end{tikzpicture}\end{center}



已知
\[
f(x)=\frac{x^2+ax+1}{x^2+1},\qquad a\in\mathbb R.
\]
因为$x^2+1>0$，定义域是$\mathbb R$，已经对称。于是
\[
f(-x)=\frac{x^2-ax+1}{x^2+1}.
\]
要使$f(-x)=f(x)$对\textbf{每个}$x$成立，必须$-ax=ax$，即$2ax=0$对每个$x$成立。取$x=1$就知道$a=0$，且代回后确实是偶函数。

\textbf{答案：}$\boxed{a=0}$。本题也不可能是奇函数，因为$f(0)=1$，而在$0$有定义的奇函数必须满足$f(0)=0$。

\medskip
\end{examplebox}
\clearpage
\begin{examplebox}
\textbf{例7｜把一个函数拆成奇、偶两部分}
\begin{center}\begin{tikzpicture}[font=\small,>=stealth,line join=round]
\draw[gray!40,rounded corners=3pt] (0,0) rectangle (14.2,10.2);
\node at (7.1,9.6) {奇偶分解：两边的平均值给偶部，偏差给奇部};
\begin{scope}[shift={(3.6,3.7)},x=1.9cm,y=0.55cm]\draw[->,gray] (-1.4,0)--(1.5999999999999999,0) node[right]{$x$};\draw[->,gray] (0,-4)--(0,8.5) node[left]{$y$};\begin{scope}\clip (-1.4,-4) rectangle (1.4,8.3);\draw[blue,thick,domain=-1.4:1.4,samples=120] plot(\x,{\x*\x+3*\x+2});\draw[teal!80!black,thick,domain=-1.4:1.4,samples=120] plot(\x,{\x*\x+2});\draw[orange!85!black,thick,domain=-1.4:1.4,samples=120] plot(\x,{3*\x});\end{scope}\draw[gray,dotted] (-1,0)--(-1,3)--(1,3)--(1,6);\fill[blue] (-1,0) circle(2pt) node[below left]{$(-1,0)$};\fill[blue] (1,6) circle(2pt) node[above left]{$(1,6)$};\fill[teal!80!black] (-1,3) circle(2pt);\fill[teal!80!black] (1,3) circle(2pt);\foreach \v in {-3,3,6}{\draw[gray] (-.05,\v)--(.05,\v) node[left]{$\v$};}\end{scope}\node[align=left,anchor=west] at (7.6,5.5) {\textcolor{blue}{$f(x)=x^2+3x+2$}\\[12pt]\textcolor{teal!80!black}{$A(x)=x^2+2$：偶部}\\[12pt]\textcolor{orange!85!black}{$B(x)=3x$：奇部}\\[15pt]$A(1)=(6+0)/2=3$\\[12pt]$B(1)=(6-0)/2=3$\\[12pt]$B(-1)=-3$\\[15pt]每个 $x$ 都有 $f(x)=A(x)+B(x)$};
\node[text width=12.8cm,align=center] at (7.1,.65) {同一坐标系中画三条真实曲线；纵轴已标实际数值，横纵比例不同。};
\end{tikzpicture}\end{center}



已知$f(x)=x^2+3x+2$，求它的偶的部分与奇的部分。
\[
\begin{aligned}
f(-x)&=x^2-3x+2,\\
\frac{f(x)+f(-x)}2&=x^2+2,\\
\frac{f(x)-f(-x)}2&=3x.
\end{aligned}
\]
\textbf{结果：}$f(x)=\underbrace{(x^2+2)}_{\text{偶函数}}+\underbrace{3x}_{\text{奇函数}}$。
\end{examplebox}
